
\documentclass{article}
\usepackage{colortbl}
\usepackage{makecell}
\usepackage{multirow}
\usepackage{supertabular}

\begin{document}

\newcounter{utterance}

\centering \large Interaction Transcript for game `eqbench', experiment `eq\_bench', episode 160 with reflexion{-}llama{-}3.1{-}8b{-}instruct.
\vspace{24pt}

{ \footnotesize  \setcounter{utterance}{1}
\setlength{\tabcolsep}{0pt}
\begin{supertabular}{c@{$\;$}|p{.15\linewidth}@{}p{.15\linewidth}p{.15\linewidth}p{.15\linewidth}p{.15\linewidth}p{.15\linewidth}}
    \# & \multicolumn{2}{c}{Player} && \multicolumn{2}{c}{Game Master} \\
    \hline

    \theutterance \stepcounter{utterance}  
    & & & \multicolumn{4}{p{0.6\linewidth}}{
        \cellcolor[rgb]{0.9,0.9,0.9}{
            \makecell[{{p{\linewidth}}}]{
                \texttt{\tiny{[P1$\langle$GM]}}
                \texttt{You have attempted this same task before and did not succeed. These are the plans you wrote after those attempts.} \\
\texttt{Trial \#1: New plan: I got stuck in a loop of analyzing the emotions expressed in the dialogue without considering the underlying dynamics of the conversation. I focused on individual statements and reactions without taking into account the overall tone and progression of the conversation. Next time, I will pay closer attention to the power dynamics and the flow of the conversation, considering how each statement builds upon the previous one and how the emotions expressed by the characters are influenced by their past interactions and relationships. I will also try to identify the underlying themes and conflicts that drive the conversation, rather than just focusing on the surface{-}level emotions expressed.} \\
\\ 
\texttt{Trial \#2: New plan: I got stuck in a loop of analyzing the emotions expressed in the dialogue without considering the underlying dynamics of the conversation. I focused on individual statements and reactions without taking into account the overall tone and progression of the conversation. Next time, I will pay closer attention to the power dynamics and the flow of the conversation, considering how each statement builds upon the previous one and how the emotions expressed by the characters are influenced by their past interactions and relationships. I will also try to identify the underlying themes and conflicts that drive the conversation, rather than just focusing on the surface{-}level emotions expressed.} \\
\\ 
\texttt{To achieve this, I will:} \\
\\ 
\texttt{1. Identify the power dynamics between Richie and Jules, considering who has the upper hand in the conversation and how this affects their emotional responses.} \\
\texttt{2. Analyze the flow of the conversation, looking for patterns and themes that emerge over time.} \\
\texttt{3. Consider the past interactions and relationships between Richie and Jules, and how these have shaped their emotional responses to each other.} \\
\texttt{4. Identify the underlying conflicts and themes that drive the conversation, and how these influence the emotions expressed by the characters.} \\
\\ 
\texttt{By taking a more nuanced and dynamic approach to analyzing the emotions expressed in the dialogue, I hope to better understand the underlying forces at play and make more accurate predictions about the emotional responses of the characters.} \\
\\ 
\texttt{Trial \#3: New plan: I got stuck in a loop of analyzing the emotions expressed in the dialogue without considering the underlying dynamics of the conversation. I focused on individual statements and reactions without taking into account the overall tone and progression of the conversation. Next time, I will pay closer attention to the power dynamics and the flow of the conversation, considering how each statement builds upon the previous one and how the emotions expressed by the characters are influenced by their past interactions and relationships. I will also try to identify the underlying themes and conflicts that drive the conversation, rather than just focusing on the surface{-}level emotions expressed.} \\
\\ 
\texttt{To achieve this, I will:} \\
\\ 
\texttt{1. Identify the power dynamics between Richie and Jules, considering who has the upper hand in the conversation and how this affects their emotional responses.} \\
\texttt{2. Analyze the flow of the conversation, looking for patterns and themes that emerge over time.} \\
\texttt{3. Consider the past interactions and relationships between Richie and Jules, and how these have shaped their emotional responses to each other.} \\
\texttt{4. Identify the underlying conflicts and themes that drive the conversation, and how these influence the emotions expressed by the characters.} \\
\\ 
\texttt{Additionally, I will also pay closer attention to the language and tone used by each character, as well as the body language and nonverbal cues that may be implied. This will help me to better understand the emotional nuances of the conversation and make more accurate predictions about the emotional responses of the characters.} \\
\\ 
\texttt{To take it a step further, I will also consider the following:} \\
\\ 
\texttt{* How does Richie's past behavior and actions influence Jules's emotional response to him?} \\
\texttt{* How does Jules's past behavior and actions influence Richie's emotional response to him?} \\
\texttt{* Are there any underlying issues or conflicts that are not explicitly stated in the dialogue, but are still influencing the emotional responses of the characters?} \\
\\ 
\texttt{By taking a more comprehensive and nuanced approach to analyzing the emotions expressed in the dialogue, I hope to better understand the underlying forces at play and make more accurate predictions about the emotional responses of the characters.} \\
\\ 
\\ 
\\ 
\texttt{Your task is to predict the likely emotional responses of a character in this dialogue:} \\
\\ 
\texttt{Richie: Jules Baxter, as I live and breathe. Never thought I'd see you at a festival again.} \\
\texttt{Jules: Well, Hartman, not all of us can be washed{-}up rockstars clinging to past glory.} \\
\texttt{Richie: Ouch. That's a bit harsh, don't you think?} \\
\texttt{Jules: You always did have a knack for understatement.} \\
\texttt{Richie: And you always had a knack for being a buzzkill. } \\
\texttt{Jules: I was the only one who ever took the music seriously, Richie. You were too busy playing the rockstar.} \\
\texttt{Richie: And you were too busy playing the martyr. } \\
\texttt{Jules: I was trying to create something meaningful. You just wanted the fame.} \\
\texttt{Richie: And you just wanted to control everything. That's why you left, isn't it?} \\
\texttt{Jules: I left because you wouldn’t listen. You never did.} \\
\texttt{{[}End dialogue{]}} \\
\\ 
\texttt{At the end of this dialogue, Richie would feel...} \\
\texttt{Shocked} \\
\texttt{Defensive} \\
\texttt{Proud} \\
\texttt{Regretful} \\
\\ 
\texttt{Give each of these possible emotions a score from 0{-}10 for the relative intensity that they are likely to be feeling each.} \\
\\ 
\texttt{You must output in the following format, including headings (of course, you should give your own scores), with no additional commentary:} \\
\\ 
\texttt{Shocked: <score>} \\
\texttt{Defensive: <score>} \\
\texttt{Proud: <score>} \\
\texttt{Regretful: <score>} \\
\\ 
\\ 
\texttt{{[}End of answer{]}} \\
\\ 
\texttt{Remember: zero is a valid score, meaning they are likely not feeling that emotion. You must score at least one emotion > 0.} \\
\\ 
\texttt{Your answer:} \\
            }
        }
    }
     \\ \\

    \theutterance \stepcounter{utterance}  
    & \multicolumn{4}{p{0.6\linewidth}}{
        \cellcolor[rgb]{0.561,0.737,0.561}{
            \makecell[{{p{\linewidth}}}]{
                \texttt{\tiny{[P1$\rangle$GM]}}
                \texttt{Shocked: 0} \\
\texttt{Defensive: 8} \\
\texttt{Proud: 0} \\
\texttt{Regretful: 9} \\
            }
        }
    }
    & & \\ \\

    \theutterance \stepcounter{utterance}  
    & & & \multicolumn{2}{p{0.3\linewidth}}{
        \cellcolor[rgb]{0.9,0.9,0.9}{
            \makecell[{{p{\linewidth}}}]{
                \texttt{\tiny{[GM$|$GM]}}
                \texttt{Shocked: 0} \\
\texttt{Defensive: 8} \\
\texttt{Proud: 0} \\
\texttt{Regretful: 9} \\
            }
        }
    }
    & & \\ \\

    \theutterance \stepcounter{utterance}  
    & & & \multicolumn{2}{p{0.3\linewidth}}{
        \cellcolor[rgb]{0.9,0.9,0.9}{
            \makecell[{{p{\linewidth}}}]{
                \texttt{\tiny{[GM$|$GM]}}
                \texttt{Shocked: 2} \\
\texttt{Defensive: 8} \\
\texttt{Proud: 4} \\
\texttt{Regretful: 6} \\
            }
        }
    }
    & & \\ \\

    \theutterance \stepcounter{utterance}  
    & & & \multicolumn{2}{p{0.3\linewidth}}{
        \cellcolor[rgb]{0.9,0.9,0.9}{
            \makecell[{{p{\linewidth}}}]{
                \texttt{\tiny{[GM$|$GM]}}
                \texttt{game\_result = LOSE} \\
            }
        }
    }
    & & \\ \\

    \theutterance \stepcounter{utterance}  
    & & \multicolumn{4}{p{0.6\linewidth}}{
        \cellcolor[rgb]{0.95,0.95,0.95}{
            \makecell[{{p{\linewidth}}}]{
                \texttt{\tiny{[P1$\leftrightarrow$P]}}
                \texttt{Post{-}episode Reflexion (not sent to GM):} \\
\texttt{New plan: I got stuck in a loop of analyzing the emotions expressed in the dialogue without considering the underlying dynamics of the conversation. I focused on individual statements and reactions without taking into account the overall tone and progression of the conversation. Next time, I will pay closer attention to the power dynamics and the flow of the conversation, considering how each statement builds upon the previous one and how the emotions expressed by the characters are influenced by their past interactions and relationships. I will also try to identify the underlying themes and conflicts that drive the conversation, rather than just focusing on the surface{-}level emotions expressed.} \\
\\ 
\texttt{To achieve this, I will:} \\
\\ 
\texttt{1. Identify the power dynamics between Richie and Jules, considering who has the upper hand in the conversation and how this affects their emotional responses.} \\
\texttt{2. Analyze the flow of the conversation, looking for patterns and themes that emerge over time.} \\
\texttt{3. Consider the past interactions and relationships between Richie and Jules, and how these have shaped their emotional responses to each other.} \\
\texttt{4. Identify the underlying conflicts and themes that drive the conversation, and how these influence the emotions expressed by the characters.} \\
\\ 
\texttt{Additionally, I will also pay closer attention to the language and tone used by each character, as well as the body language and nonverbal cues that may be implied. This will help me to better understand the emotional nuances of the conversation and make more accurate predictions about the emotional responses of the characters.} \\
\\ 
\texttt{To take it a step further, I will also consider the following:} \\
\\ 
\texttt{* How does Richie's past behavior and actions influence Jules's emotional response to him?} \\
\texttt{* How does Jules's past behavior and actions influence Richie's emotional response to him?} \\
\texttt{* Are there any underlying issues or conflicts that are not explicitly stated in the dialogue, but are still influencing the emotional responses of the characters?} \\
\\ 
\texttt{However, I realize that I still got stuck in a loop of similar analysis. To break this loop, I will take a step back and re{-}evaluate the conversation from a different perspective. I will ask myself:} \\
\\ 
\texttt{* What is the underlying message or theme that Richie and Jules are trying to convey to each other?} \\
\texttt{* How do their past interactions and relationships influence their current conversation?} \\
\texttt{* What are the underlying conflicts or issues that are driving the conversation?} \\
\\ 
\texttt{By taking a more nuanced and dynamic approach to analyzing the emotions expressed in the dialogue, I hope to better understand the underlying forces at play and make more accurate predictions about the emotional responses of the characters.} \\
            }
        }
    }
    & \\ \\

\end{supertabular}
}

\end{document}
